\documentclass[11pt]{article}

\usepackage[margin=1.1in]{geometry}
\usepackage{amsmath,amssymb,amsthm}
\usepackage{booktabs}
\usepackage{hyperref}
\usepackage{enumitem}
\usepackage{microtype}
\usepackage{listings}
\usepackage{xcolor}

\lstset{
  basicstyle=\small\ttfamily,
  keywordstyle=\color{blue!70!black},
  commentstyle=\color{green!50!black},
  breaklines=true,
  frame=single,
  xleftmargin=1em,
  framexleftmargin=0.5em,
  columns=fullflexible,
  % Code shown with ASCII syntax; actual source uses Unicode equivalents.
}

\newtheorem{theorem}{Theorem}
\newtheorem{lemma}[theorem]{Lemma}
\newtheorem{definition}[theorem]{Definition}
\newtheorem{proposition}[theorem]{Proposition}
\newtheorem{remark}[theorem]{Remark}

\newcommand{\Bool}{\{0,1\}}
\newcommand{\Sn}{S_n}
\newcommand{\Ztwo}{\mathbb{Z}_2}
\newcommand{\CC}{\mathsf{C}}

\title{Empirical Circuit Complexity:\\
  Formal Proofs for Small Boolean Functions}
\author{Samuel Schlesinger}
\date{}

\begin{document}
\maketitle

\begin{abstract}
We describe a research program to determine the exact circuit complexity of
every Boolean function $\Bool^n \to \Bool$ for small~$n$, with all results
formally verified in Lean~4. Our circuit model uses fan-in~2 AND gates with
free negation on wires. We exploit NPN equivalence---input permutation, input
negation, output negation---to reduce the proof burden by orders of magnitude:
from 65{,}536 functions at $n=4$ down to 222 representatives. The proofs are
written by AI within a trust architecture where Lean's type checker is the sole
arbiter of correctness.
\end{abstract}

\section{Introduction}

Circuit complexity---the minimum number of gates needed to compute a Boolean
function---is among the most fundamental measures in computational complexity
theory. Despite decades of effort, proving superlinear circuit lower bounds for
explicit functions remains one of the great open
problems~\cite{razborov, arora-barak}.

The asymptotic theory has produced deep structural results: Lupanov's upper
bound showing almost all functions require $\sim 2^n/n$
gates~\cite{lupanov}, and the Shannon counting argument~\cite{shannon} showing
most functions are hard. For small~$n$, exact circuit complexities have been
determined computationally, notably in Knuth's exhaustive
treatment~\cite{knuth-4a}. But these results are computational artifacts, not
proofs: they are outputs of programs that could, in principle, contain bugs.

This paper describes a program to produce \emph{formal proofs} of the exact
circuit complexity of every Boolean function on a small number of inputs. For
$n$ input bits, there are $2^{2^n}$ Boolean functions. The circuit complexity
of each is a natural number, and we prove it correct---both by exhibiting an
optimal circuit and by ruling out all smaller ones---as a machine-checked
theorem in Lean~4.

Three aspects are worth noting:

\begin{enumerate}[leftmargin=*]
\item \textbf{Formal verification.} Every result is a theorem that Lean's type
  checker has accepted. If any claim is wrong, the project does not compile.

\item \textbf{Symmetry exploitation.} We formally prove that circuit
  complexity is invariant under the NPN equivalence group, reducing the number
  of independent proofs by up to $295\times$ at $n=4$.

\item \textbf{AI-assisted proof generation.} The proof terms are written by
  large language models within a trust architecture that confines the AI to
  supplying evidence. Incorrect proofs are rejected by the type checker. This
  enables scale without sacrificing rigor.
\end{enumerate}

Our initial target is $n = 1$ through $n = 4$: 242 representative proofs
covering all 65{,}812 Boolean functions. Beyond the database itself, we are
interested in the structural questions it surfaces: what is the distribution of
circuit complexities? How do optimal circuits and lower bound arguments change
as $n$ grows?

The source code is available at
\url{https://github.com/sschlesinger/empirical-circuit-complexity}.\footnote{%
Repository URL subject to change; the project may move or be renamed.}

\section{Circuit Model}

\begin{definition}
A \emph{Boolean circuit} with $n$ inputs and $k$ gates is a sequence of $k$
fan-in~2 AND gates in topological order, together with a designated output
wire. Each wire reference (gate input or circuit output) consists of a wire
index and a \emph{negation bit}; when the negation bit is set, the signal on
the wire is inverted. The \emph{size} of the circuit is $k$, the number of AND
gates.
\end{definition}

\begin{definition}
The \emph{circuit complexity} of a Boolean function $f : \Bool^n \to \Bool$,
written $\CC(f)$, is the minimum $k$ such that some circuit of size $k$
computes~$f$.
\end{definition}

Formally, a circuit is a pair $(\mathbf{g}, r)$ where $\mathbf{g}$ is a
\texttt{GateList}---an inductively defined sequence of gates, each referencing
only earlier wires---and $r$ is a \texttt{Ref} (a wire index with negation bit)
designating the output. Gate~$i$ can reference any of the $n + i$ available
wires: $n$ inputs plus $i$ preceding gate outputs.

The formalization is self-contained: it uses Lean~4 with no external
dependencies (no Mathlib or other libraries). This is a deliberate choice to
keep the trusted base minimal, though we may reconsider as the proofs grow more
complex.

\paragraph{Why this basis.} The AND/NOT basis is functionally complete.
Treating NOT as a zero-cost wire annotation reflects the standard convention
that negation is ``free''~\cite{arora-barak}. This model is equivalent to
counting \{AND, OR\} gates when both are available, since
$\text{OR}(a,b) = \lnot\text{AND}(\lnot a, \lnot b)$ by De~Morgan's law.

\paragraph{Constants cost one gate.} There are no free constant wires. The
cheapest way to produce a constant output is $\text{AND}(x, \lnot x) =
\mathsf{false}$ (one gate), optionally negated for $\mathsf{true}$.

\section{Trust Architecture}

A central design goal is a sharp trust boundary between \emph{what is claimed}
and \emph{who generates the evidence}. For each input size $n$, three kinds of
files are produced:

\begin{enumerate}[leftmargin=*]
\item \textbf{Definitions} (\texttt{Defs}): Generated by a human-reviewed
  Python program. Each Boolean function is encoded as a Lean~4 function
  \texttt{(Fin $n$ → Bool) → Bool} via a nested if-then-else tree matching
  its truth table.

\item \textbf{Proofs}: Written by AI. Contains proof terms: upper bounds
  (exhibiting circuits) and lower bounds (ruling out smaller circuits). These
  are the only files not produced by the trusted generator.

\item \textbf{Statements}: Generated by the same trusted program. Imports
  both Definitions and Proofs and states the circuit complexity theorems,
  drawing proof terms from the Proofs module.
\end{enumerate}

The trust model is: we trust the generator (which we review) and Lean's type
checker. The AI is \emph{untrusted}---it may hallucinate circuits, miscalculate
complexities, or produce invalid proof terms. Any such error causes the
Statements file to fail type-checking. Lean's module system enforces this
boundary: the Statements file imports the Proofs file and requires that the
exported terms have the correct types.

For each representative function, the proof file must export:
\begin{itemize}[leftmargin=*]
\item \texttt{f\_\{idx\}\_size : Nat} --- the claimed circuit complexity $k$,
\item \texttt{f\_\{idx\}\_upper : HasCircuitOfSize f $k$} --- a
  witness circuit of size $k$,
\item \texttt{f\_\{idx\}\_lower : $\forall$ j, j < $k$ $\to$ $\lnot$ HasCircuitOfSize f $j$}
  --- a proof that no smaller circuit suffices.
\end{itemize}

\section{Symmetry Exploitation}
\label{sec:symmetry}

Several transformations on circuits are \emph{free}: they change which function
is computed without changing the number of gates. We identify three such
transformations and formally prove that each preserves circuit complexity.

\subsection{Three free transformations}

Write $f^\sigma$ for the function $x \mapsto f(x_{\sigma(0)}, \ldots,
x_{\sigma(n-1)})$, i.e., $f$ with inputs permuted by $\sigma$.
Write $f^{\oplus i}$ for the function obtained by negating input~$i$:
$x \mapsto f(x_0, \ldots, \lnot x_i, \ldots, x_{n-1})$.

\begin{theorem}[Input permutation invariance]
For any $\sigma \in \Sn$,\; $\CC(f^\sigma) = \CC(f)$.
\end{theorem}
\noindent\textit{Proof idea.} Relabel each input wire reference $i \mapsto
\sigma(i)$. No gates are added or removed.

\begin{theorem}[Input negation invariance]
For any input index $i$,\; $\CC(f^{\oplus i}) = \CC(f)$.
\end{theorem}
\noindent\textit{Proof idea.} Flip the negation bit on every reference to input
wire $i$. No gates are added or removed.

\begin{theorem}[Output negation invariance]
$\CC(\lnot f) = \CC(f)$.
\end{theorem}
\noindent\textit{Proof idea.} Flip the negation bit on the output wire. No
gates are added or removed.

\medskip
All three theorems and their converses are stated and proved in Lean~4 in our
\texttt{Circuits/Basic.lean} module. The output negation proof is complete. The
input permutation and input negation proofs depend on evaluation lemmas
(showing that the modified circuit evaluates to the expected function) whose
inductive proofs over the dependent gate list type are stated but deferred; the
complexity invariance theorems are proved assuming these lemmas.

\subsection{NPN equivalence}

These three operations generate a group known in logic synthesis as the
\emph{NPN equivalence group}~\cite{npn-corrected}:
\[
  G = (\Sn \ltimes \Ztwo^n) \times \Ztwo,
\]
of order $2^{n+1} \cdot n!$, where $\Sn$ acts on $\Ztwo^n$ by permuting
coordinates. Two functions are NPN-equivalent if and only if one can be obtained
from the other by permuting inputs, negating inputs, and negating the output.
The NPN class counts form OEIS sequence A000370~\cite{oeis-a000370}.

\begin{remark}
NPN equivalence is the maximal equivalence relation on Boolean functions
preservable by free wire operations in this model. The argument is
straightforward: the only operations that preserve gate count are changes to
wire indices and negation bits. On input wires, index changes yield $\Sn$ and
negation bit changes yield $\Ztwo^n$; on the output wire, negation bit changes
yield the final $\Ztwo$ factor. Any other input transformation---for instance,
replacing $x_i$ with $x_i \oplus x_j$---requires computing a new function of
two wires, which costs at least one gate.
\end{remark}

\subsection{Counting representatives via Burnside's lemma}

The number of NPN equivalence classes is computed by Burnside's lemma:
\[
  |X / G| = \frac{1}{|G|} \sum_{g \in G} |\mathrm{Fix}(g)|.
\]
Each group element $g = (\sigma, \varepsilon, \delta)$ induces a permutation
$\pi$ on the $2^n$ input indices: $\pi(j)$ is obtained by permuting the bits of
$j$ according to $\sigma$ and then XORing with $\varepsilon$. A truth table is
fixed by $g$ if and only if $f(j) = f(\pi(j)) \oplus \delta$ for every input
index~$j$.

When $\delta = 0$, the truth table must be constant on each cycle of $\pi$,
giving $2^{c(\pi)}$ fixed points where $c(\pi)$ is the number of cycles.
When $\delta = 1$, values must alternate around each cycle:
$f(j) \neq f(\pi(j))$. Following this constraint around a cycle of length~$L$,
we find $f(j) = f(\pi^L(j)) \oplus (L \bmod 2) = f(j) \oplus (L \bmod 2)$.
If $L$ is odd this forces $f(j) \neq f(j)$, a contradiction, so all cycles
must have even length; when they do, there are $2^{c(\pi)}$ fixed points.

\begin{table}[ht]
\centering
\begin{tabular}{@{}rrrrr@{}}
\toprule
$n$ & Functions & $|G|$ & Representatives & Reduction \\
\midrule
1 & 4 & 4 & 2 & $2\times$ \\
2 & 16 & 16 & 4 & $4\times$ \\
3 & 256 & 96 & 14 & $18\times$ \\
4 & 65{,}536 & 768 & 222 & $295\times$ \\
5 & ${\sim}4.3 \times 10^9$ & 7{,}680 & 616{,}126 & $6{,}971\times$ \\
6 & ${\sim}1.8 \times 10^{19}$ & 92{,}160 & ${\sim}2.0 \times 10^{14}$ & $92{,}117\times$ \\
7 & ${\sim}3.4 \times 10^{38}$ & 1{,}290{,}240 & ${\sim}2.6 \times 10^{32}$ & $1{,}290{,}240\times$ \\
\bottomrule
\end{tabular}
\caption{NPN equivalence class counts. For large $n$, the reduction factor
  approaches $|G|$, indicating almost all functions have trivial stabilizer.}
\label{tab:reps}
\end{table}

At $n = 4$, proving complexity for 222 representatives covers all 65{,}536
functions. For $n \le 4$, only $2 + 4 + 14 + 222 = 242$ independent proofs are
needed.

\section{Proof Techniques}

Two kinds of proofs are required for each function: an upper bound (exhibit a
circuit) and a lower bound (rule out all smaller circuits).

\subsection{Upper bounds}

An upper bound constructs a concrete circuit and proves it computes the target
function. The proof obligation is
$\forall\, x,\; \texttt{eval}(c, x) = f(x)$.
Because both the circuit and the function definition are fully concrete, Lean's
\texttt{simp} tactic---given the evaluation definitions---reduces both sides to
identical Boolean expressions and closes the goal automatically.

\subsection{Lower bounds}

Lower bound proofs show that no circuit of size less than $k$ computes~$f$.
This requires universal quantification over all circuits of each smaller size.
Several patterns emerge at small~$n$:

\paragraph{Vacuous.} When $\CC(f) = 0$, the lower bound
$\forall\, j < 0, \ldots$ is vacuously true.

\paragraph{Ruling out size 0.} A size-0 circuit is just a (possibly negated)
input wire. A general lemma handles all constant functions: evaluating on the
all-true and all-false inputs, a size-0 circuit must return different values,
but a constant function returns the same value. Contradiction. For non-constant,
non-literal functions, the proof case-splits on which input wire and negation
bit, providing a counterexample input for each case.

\paragraph{Ruling out size $\ge 1$.} The proof must enumerate all possible
gate configurations. For a single gate, there are $(2n)^2$ choices for the gate
inputs (each of $n$ wires with a negation bit, two inputs) and $2(n+1)$
choices for the output reference. The proof destructures the circuit,
case-splits on all choices, and provides counterexample inputs. For size~$\ge 2$,
the enumeration grows combinatorially.

\subsection{Example: constant false on one input}

To illustrate the proof structure, consider $f_0$ on $n = 1$: the
constant-false function. Its complexity is~1.

\paragraph{Upper bound.} We exhibit the circuit
$\text{AND}(x_0, \lnot x_0)$ with output referencing the gate (non-negated):

\begin{lstlisting}
def f_0_upper : HasCircuitOfSize f_0 1 :=
  ({gates := .cons .nil {lhs := {index := (0 : Fin 2),
                                  negated := false},
                          rhs := {index := (0 : Fin 2),
                                  negated := true}},
    output := {index := (1 : Fin 2), negated := false}},
   fun input => by
     simp [Circuit.eval, Ref.eval, GateList.eval,
           Gate.eval, extendEnv, f_0])
\end{lstlisting}

The anonymous constructor provides the circuit term; \texttt{simp} verifies
that evaluating it on any input produces $\mathsf{false}$, matching $f_0$.

\paragraph{Lower bound.} We must show no size-0 circuit computes $f_0$. Since
$f_0(\mathbf{1}) = f_0(\mathbf{0}) = \mathsf{false}$, the general constant
lemma applies directly:

\begin{lstlisting}
def f_0_lower :
    forall j, j < 1 -> Not (HasCircuitOfSize f_0 j) := by
  intro j hj
  have : j = 0 := by omega
  subst this
  exact no_size0_of_constant (by simp [f_0])
\end{lstlisting}

\noindent The \texttt{omega} tactic establishes $j = 0$; then
the general lemma rules out size~0.

\begin{remark}
A recurring implementation concern: Lean's \texttt{simp} tactic diverges when
applied to circuit evaluation terms containing free Boolean variables (e.g., an
unresolved negation bit). The workaround is to case-split on all Boolean
parameters \emph{before} invoking \texttt{simp}.
\end{remark}

\section{Results}

\subsection{Completed: $n = 1$}

All 4 functions (2 NPN representatives) have been formally verified:

\begin{center}
\begin{tabular}{@{}clcl@{}}
\toprule
Index & Function & $\CC$ & Witness \\
\midrule
0 & constant $0$ & 1 & $\text{AND}(x_0, \lnot x_0)$ \\
1 & $\lnot x_0$ & 0 & wire with negation \\
\bottomrule
\end{tabular}
\end{center}

The two non-representatives are $x_0$ (output negation of $\lnot x_0$) and
constant~$1$ (output negation of constant~$0$).

\subsection{Expected: $n = 2$}

The 16 functions reduce to 4 NPN representatives. Formal proofs are in
progress; the expected complexities (matching Knuth's computed
values~\cite{knuth-4a}) are:

\begin{center}
\begin{tabular}{@{}clcl@{}}
\toprule
Index & Function & $\CC$ (expected) & Witness \\
\midrule
0 & constant $0$ & 1 & $\text{AND}(x_0, \lnot x_0)$ \\
1 & $\lnot x_0 \land \lnot x_1$ & 1 & $\text{AND}(\lnot x_0, \lnot x_1)$ \\
3 & $\lnot x_1$ & 0 & wire with negation \\
6 & $x_0 \oplus x_1$ & 3 & see below \\
\bottomrule
\end{tabular}
\end{center}

The XOR function is notable: it requires 3 AND gates, making it the most
expensive 2-input function. One optimal circuit computes it as
$\lnot(\lnot(x_0 \land \lnot x_1) \land \lnot(\lnot x_0 \land x_1))$:
the two inner gates compute the ``half-XORs'' $x_0 \land \lnot x_1$ and
$\lnot x_0 \land x_1$, and the outer gate takes their NOR and negates.
The lower bound proof---ruling out all circuits of sizes 0, 1, and 2---will
require significant case analysis.

\subsection{Baseline observations}

At $n = 1$, the results establish baseline behavior for the cost model:

\begin{itemize}[leftmargin=*]
\item \textbf{Negation is free in practice.} The literal $\lnot x_0$ has
  complexity~0. This confirms the model: NOT is a wire annotation, not a gate.

\item \textbf{Constants are not free.} With no constant wires available,
  producing a constant costs one gate ($x \land \lnot x$). This is a
  consequence of the model choice and sets the floor for all non-literal
  functions.

\item \textbf{NPN reduction is effective.} Even at $n = 1$, the four functions
  collapse to two classes. At $n = 2$, 16 functions collapse to 4 classes. At
  $n = 4$, the $295\times$ reduction makes the project feasible.
\end{itemize}

\noindent These are not surprises---they follow directly from the
definitions---but they validate that the formal model behaves as intended before
we scale to larger~$n$ where the results are no longer obvious.

\section{Related Work}

\paragraph{Exact circuit complexity.}
Knuth~\cite{knuth-4a} gives an extensive computational treatment of circuit
complexity for small Boolean functions, including exhaustive search for
optimal circuits and tables of exact values. Our work produces the same
values but as formally verified theorems rather than computational outputs. The
formal proofs also serve a different purpose: they are explicit, inspectable
objects from which proof strategies can be extracted and analyzed.

\paragraph{NPN classification.}
NPN equivalence is a standard tool in logic synthesis for reducing the space of
Boolean functions~\cite{npn-corrected}. The class counts (OEIS
A000370~\cite{oeis-a000370}) have been tabulated for $n$ up to at least~6. We
apply this classification not for synthesis but to reduce the number of formal
proofs required.

\paragraph{Formal verification of combinatorics.}
The use of proof assistants for large combinatorial results has precedent: the
Flyspeck project~\cite{flyspeck} (Kepler conjecture), the Four Color
Theorem formalization~\cite{four-color}, and various results in Lean's
Mathlib~\cite{mathlib}. Our work is distinguished by its focus on circuit
complexity and its use of AI-generated proof terms within a minimal trusted base.

\paragraph{AI-assisted theorem proving.}
Recent work has shown the potential of large language models for formal proof
generation~\cite{llm-lean}. Our trust architecture is designed to exploit this:
the AI generates candidate proof terms, and the type checker accepts or rejects
them. The human role is limited to reviewing the trusted generator and
evaluating the proof strategies that emerge.

\section{Open Questions}

The systematic, verified nature of this effort surfaces several concrete
questions that we intend to investigate as the database grows.

\paragraph{Complexity distribution.} What is the distribution of $\CC(f)$ over
NPN representatives at each~$n$? How does the maximum complexity grow? The
Lupanov bound gives $\max_f \CC(f) \le (1 + o(1)) \cdot 2^n / n$
asymptotically~\cite{lupanov}; our data will provide exact values for
small~$n$.

\paragraph{Structure of hard functions.} Which NPN classes have the highest
complexity? Do they share structural features? XOR is the hardest function at
$n = 2$; what are its counterparts at $n = 3$ and $n = 4$?

\paragraph{Lower bound composition.} Can lower bound proofs be composed
inductively across levels of~$n$? Shannon decomposition---fixing one variable
and reducing to two $(n{-}1)$-input subfunctions---is a natural strategy. Our
formal framework supports cross-size imports, enabling proofs at level~$n$ to
cite verified results from level~$n - 1$.

\paragraph{Automation frontier.} At $n = 4$ (222 representatives), AI-assisted
proof generation is feasible with per-function interaction. At $n = 5$
(616{,}126 representatives), it requires genuine automation: batch proof search
guided by patterns from smaller~$n$. Upper bounds are likely automatable (search
for circuits); lower bounds are the bottleneck.

\paragraph{Extracting proof structure.} Every formal proof is an explicit term
available for inspection. The counterexample inputs used in lower bound proofs,
the circuit topologies that appear in optimal solutions, and the case-analysis
patterns that close goals are all concrete data. Systematically mining this
data may reveal structure invisible to the asymptotic theory.

\section{Conclusion}

We have described a research program to determine the exact circuit complexity
of all Boolean functions for small~$n$, with formal verification in Lean~4 as
the standard of correctness. The NPN symmetry group reduces the proof burden by
orders of magnitude, from 65{,}812 functions to 242 representatives for $n \le
4$. An AI-assisted workflow generates proof candidates that the type checker
verifies independently, enabling scale without sacrificing rigor.

The $n = 1$ case is complete; $n = 2$ through $n = 4$ are in progress. As the
verified database grows, we expect it to serve as a testbed for questions about
circuit complexity---one where every claimed fact is machine-checked and every
proof term is available for inspection.

\begin{thebibliography}{99}

\bibitem{razborov}
  A.~A.~Razborov, S.~Rudich,
  ``Natural proofs,''
  \textit{J.~Comput.~Syst.~Sci.}, 55(1):24--35, 1997.

\bibitem{arora-barak}
  S.~Arora, B.~Barak,
  \textit{Computational Complexity: A Modern Approach},
  Cambridge University Press, 2009.

\bibitem{lupanov}
  O.~B.~Lupanov,
  ``On the synthesis of contact networks''
  (\textit{O sinteze kontaktnykh skhem}),
  \textit{Dokl.~Akad.~Nauk SSSR}, 119(1):23--26, 1958.

\bibitem{shannon}
  C.~E.~Shannon,
  ``The synthesis of two-terminal switching circuits,''
  \textit{Bell Syst.~Tech.~J.}, 28(1):59--98, 1949.

\bibitem{knuth-4a}
  D.~E.~Knuth,
  \textit{The Art of Computer Programming}, Vol.~4A:
  \textit{Combinatorial Algorithms, Part 1},
  Addison-Wesley, 2011.
  See Section~7.1.2.

\bibitem{npn-corrected}
  D.~Debnath, T.~Sasao,
  ``Efficient computation of canonical form for Boolean matching in
  large libraries,''
  \textit{Proc.~ASP-DAC}, 2004.

\bibitem{oeis-a000370}
  OEIS Foundation,
  ``A000370: Number of NPN-equivalence classes of Boolean functions of $n$
  variables,''
  \url{https://oeis.org/A000370}.

\bibitem{flyspeck}
  T.~Hales et~al.,
  ``A formal proof of the Kepler conjecture,''
  \textit{Forum Math.~Pi}, 5:e2, 2017.

\bibitem{four-color}
  G.~Gonthier,
  ``Formal proof---the four-color theorem,''
  \textit{Notices AMS}, 55(11):1382--1393, 2008.

\bibitem{mathlib}
  The mathlib Community,
  ``The Lean mathematical library,''
  \textit{Proc.~CPP}, 2020.

\bibitem{llm-lean}
  K.~Yang et~al.,
  ``LeanDojo: Theorem proving with retrieval-augmented language models and
  data extraction from Lean 4,''
  \textit{NeurIPS}, 2023.

\end{thebibliography}

\end{document}
